\documentclass[11pt]{article}
\input{style-pc}
\usepackage{array}

\begin{document}

\entetesujet{Sujet bac 2018 -- Physique-Chimie -- Série D}

% ============================ CHIMIE ============================
\matiere{CHIMIE}{8 points}

\partie{A : vérification des connaissances}

\noindent\textbf{1. Texte à trous.} Recopie puis complète les quatre (04) mots ou groupes de mots manquants dans la phrase suivante parmi les mots ci-après : \textit{oxydant, potentiel, électron, réducteur, oxydoréduction}.

\textit{Dans une réaction d'$\cdots\cdots$, c'est l'$\cdots\cdots$ le plus fort qui $\cdots\cdots$ avec le $\cdots\cdots$ le plus fort.}

\medskip
\noindent\textbf{2. Questions à choix multiples.} Choisis la bonne réponse parmi les propositions suivantes. Exemple : $2.g = g_3$.

\noindent\textbf{2.a.} L'énergie de l'atome d'hydrogène est donnée par la relation :
\begin{enumerate}[label=a\textsubscript{\arabic*}.]
  \item $E_n = \dfrac{13{,}6}{n^2}$ (J) ;
  \item $E_n = -\dfrac{13{,}6}{n^2}$ (eV) ;
  \item $E_n = -13{,}6\,n^2$ (eV).
\end{enumerate}

\noindent\textbf{2.b.} Le niveau d'énergie le plus bas de l'atome d'hydrogène correspond :
\begin{enumerate}[label=b\textsubscript{\arabic*}.]
  \item Au premier état excité ;
  \item Au niveau d'ionisation ;
  \item À l'état fondamental.
\end{enumerate}

\noindent\textbf{2.c.} En dissolvant un soluté de masse $m$ dans un solvant de masse $m'$, l'abaissement cryométrique de cette solution diluée non électrolysable est :
\begin{enumerate}[label=c\textsubscript{\arabic*}.]
  \item $\Delta\Theta = K\dfrac{m'}{mM}$ ;
  \item $\Delta\Theta = K\dfrac{m}{m'M}$ ;
  \item $\Delta\Theta = K\dfrac{M}{m\,m'}$.
\end{enumerate}

\noindent\textbf{2.d.} Le pH d'une dibase forte en solution aqueuse est :
\begin{enumerate}[label=d\textsubscript{\arabic*}.]
  \item $-\log C_b$ ;
  \item $14 - \log C_b$ ;
  \item $14 + \log 2C_b$ ;
  \item $14 + \log C_b$.
\end{enumerate}

\medskip
\noindent\textbf{3. Question à réponse courte.} Dans l'équation de la réaction suivante : $\ch{2\,I^- + H_2O_2 + 2\,H_3O^+ \longrightarrow I_2 + 4\,H_2O}$. Établis une relation
\begin{enumerate}[label=\alph*.]
  \item Entre les vitesses de $\ch{H_2O_2}$ et de $\ch{H_2O}$.
  \item Entre les vitesses de $\ch{I^-}$ et de $\ch{I_2}$.
\end{enumerate}

\partie{B : application des connaissances}
Les différents niveaux d'énergie $E_n$ de l'atome d'hydrogène sont donnés par la formule : $E_n = -\dfrac{13{,}6}{n^2}$ avec $E_n$ en eV et $n \in \N^*$.

\begin{enumerate}
  \item Calcule en eV, $E_1$, $E_2$, $E_3$, $E_4$, $E_5$, $E_\infty$.
  \item Fais le diagramme d'énergie correspondant. Échelle : 1 cm pour 1 eV.
  \item Calcule :
  \begin{enumerate}[label=\alph*.]
    \item L'énergie minimale de la série de Balmer.
    \item L'énergie maximale de la série de Balmer.
  \end{enumerate}
  \item Déduis :
  \begin{enumerate}[label=\alph*.]
    \item La longueur d'onde $\lambda_1$ la plus courte de la série de Balmer.
    \item La longueur d'onde $\lambda_2$ la plus grande de la même série.
  \end{enumerate}
\end{enumerate}
\noindent On donne : la constante de Planck $h = 6{,}62\cdot10^{-34}$ J$\cdot$s ; la célérité $c = 3\cdot10^8$ m/s ; $1$ eV $= 1{,}6\cdot10^{-19}$ J.

% ============================ PHYSIQUE ============================
\matiere{PHYSIQUE}{12 points}

\partie{A : vérification des connaissances}

\noindent\textbf{1. Réarrangement.} La phrase suivante a été écrite en désordre. Reproduis puis ordonne-la.

\textit{La formation / prouve / de la lumière / des interférences / la nature / lumineuses / ondulatoire.}

\medskip
\noindent\textbf{2. Appariement.} Relie un élément-question de la colonne A à un élément-réponse de la colonne B qui lui correspond. Exemple : $A_5 = B_8$.

\begin{center}\renewcommand{\arraystretch}{1.8}
\begin{tabular}{|p{0.36\linewidth}|p{0.5\linewidth}|}
\hline
\textbf{Colonne A} & \textbf{Colonne B} \\
\hline
$A_1$ : différence de marche & $B_1$ : la tension en avance de phase sur l'intensité \\
\hline
$A_2$ : circuit inductif & $B_2$ : $\dfrac{\lambda D}{a}$ \\
\hline
$A_3$ : circuit capacitif & $B_3$ : la tension en retard de phase sur l'intensité \\
\hline
$A_4$ : interfrange & $B_4$ : $\dfrac{ax}{D}$ \\
\hline
\end{tabular}
\end{center}

\noindent\textbf{3. Question à réponse courte.} Définis un pendule pesant.

\partie{B : application des connaissances}
La cathode d'une cellule photoélectrique au potassium est éclairée par deux radiations lumineuses, l'une de longueur d'onde $\lambda_1 = 0{,}49$ µm, l'autre de longueur d'onde $\lambda_2 = 0{,}66$ µm. Le travail d'extraction d'un électron du potassium est $W_0 = 2{,}25$ eV.

\begin{enumerate}
  \item \begin{enumerate}[label=\alph*.]
    \item Détermine la longueur d'onde seuil.
    \item Laquelle des deux radiations provoque l'effet photoélectrique ?
  \end{enumerate}
  \item La cellule est maintenant éclairée par la radiation de longueur d'onde $\lambda_1 = 0{,}49$ µm. La puissance rayonnante reçue par la cathode est $P = 9\cdot10^{-7}$ W. L'intensité du courant de saturation dans le circuit vaut $I_s = 4\cdot10^{-8}$ A. Détermine :
  \begin{enumerate}[label=\alph*.]
    \item La vitesse maximale de sortie des électrons à la cathode.
    \item Le rendement quantique de la cellule.
  \end{enumerate}
\end{enumerate}
\noindent Données : $h = 6{,}62\cdot10^{-34}$ J$\cdot$s ; $c = 3\cdot10^8$ m/s ; $e = 1{,}6\cdot10^{-19}$ C ; $1$ µm $= 10^{-6}$ m ; $1$ eV $= 1{,}6\cdot10^{-19}$ J ; $m_e = 9{,}1\cdot10^{-31}$ kg.

\partie{C : résolution d'un problème}
À la suite d'une expérience sur la propagation des ondes le long d'une corde, on veut construire les graphes représentant l'élongation d'un point vibrant en fonction du temps $t$ et l'aspect de la corde à un instant donné. Pour cela, on considère une lame horizontale dont l'extrémité $A$ vibre verticalement. En $A$ est fixée une corde horizontale de longueur $L = AB = 1{,}20$ m, de masse $m = 24$ g soumise à une force $F$. La fréquence des vibrations de la lame est $N = 50$ Hz et les vibrations se propagent le long de la corde à la célérité $c = 20$ m/s. Un système supprime la réflexion des ondes à l'extrémité $B$ du fil.

\begin{center}
\begin{tikzpicture}[>=Stealth,scale=1]
  % vibreur
  \draw[pattern=north east lines] (-0.9,-0.12) rectangle (0,0.12);
  \node[left] at (-0.9,0){vibreur};
  \node[above] at (0.4,0.1){$A$};
  \draw (0,0)--(5.5,0);
  \draw (6,0) circle (0.5);
  \fill (6,0) circle (1pt);
  \node[above] at (6,0.55){$B$};
  % fil vertical + tension
  \draw (6.5,0)--(6.5,-1.2);
  \draw[->] (6.5,-1.55)--(6.5,-0.7) node[right]{$\vec{F}$};
  % triangle support
  \draw (6.5,-1.2)--(6,-2.1)--(7,-2.1)--cycle;
  \draw[fill=gray!30] (6.35,-1.95) rectangle (6.65,-1.75);
  % poids
  \draw[->] (6.5,-2.1)--(6.5,-2.9) node[right]{$\vec{P}$};
\end{tikzpicture}
\end{center}

\begin{enumerate}
  \item Calcule :
  \begin{enumerate}[label=\alph*.]
    \item La longueur d'onde du mouvement vibratoire.
    \item La tension $F$ du fil.
  \end{enumerate}
  \item L'extrémité $A$ de la lame a un mouvement sinusoïdal d'amplitude $a = 10$ mm.
  \begin{enumerate}[label=\alph*.]
    \item Établis l'équation horaire du mouvement de $A$. (On prendra comme origine des temps l'instant où l'élongation de $A$ est nulle et croissante.)
    \item Établis l'équation horaire d'un point $C$ du fil tel que $AC = x = 70$ cm.
  \end{enumerate}
  \item Représente l'élongation $y_C(t)$ du mouvement du point $C$.
  \item Représente l'aspect de la corde $AB$ à l'instant $t = 0{,}025$ s.
\end{enumerate}

\end{document}
